%% =====================================================================
%%  第1章 はじめに
%%  書くこと：背景 → 課題 → 目的 → 提案の概要 → 貢献 → 本論文の構成
%%  この章だけで「なぜこの研究が必要で，何をしたのか」がわかるようにする．
%% =====================================================================
\chapter{はじめに}
\label{chap:introduction}

\section{研究の背景}
\label{sec:background}

\begin{writingnote}
研究の背景を述べる．社会的・学術的な文脈から書き始め，本研究が扱う問題へと話題を絞っていく．
事実や主張には根拠を示し，必要に応じて文献を引用する．
題目に含まれる用語などの主要な用語は，初めて使うところで簡単に説明する（厳密な定義は「準備」の章で述べる）．

%% 段落を変えるときは空行を 1 行入れる．\\ で改行しないこと．
次の段落では，主な既存研究の成果に簡単に触れたうえで，まだ解決されていない課題を明らかにする．
既存研究の詳しい整理は「関連研究」の章で行う．
読者が「その課題を解決する必要がある」と納得できるように，具体例や数値を挙げて説明するとよい．

%% 論文の構成や文章の書き方は，次の参考書が役に立つ（引用の例も兼ねている）．
論文の構成や文章の書き方については，文献~\cite{kinoshita1981rikakei,sakai2015kore}が参考になる．
\end{writingnote}

近年，……において○○が課題となっている．
既存の方法では……であるため，……を検討する必要がある．

\section{研究の目的}
\label{sec:purpose}

\begin{writingnote}
本研究の目的を 1〜2 文で明確に述べる．
目的は前の節で述べた課題に対応させ，「結論」の章で達成できたかどうかを判断できるように具体的に書く．
続いて，目的を達成するための提案手法の概要と，実験で何を確かめたかを簡潔に述べる．
優位性を前提にせず，実験で確認できた性能，適用条件，限界を貢献として整理する．
この節は，ほかの章を書き終えてから見直し，目的と貢献が「提案手法」と「数値実験」の章で示した内容と一致しているかを確かめる．
\end{writingnote}

本研究の目的は，……を明らかにすることである．
本研究の貢献は次のとおりである．
\begin{itemize}
  \item ○○問題を△△として定式化した．
  \item ○○を用いた△△問題の解法を提案した．
  \item 数値実験により，提案手法の性能，適用条件および限界を明らかにした．
\end{itemize}

\section{本論文の構成}
\label{sec:organization}

%% 章・節・図・表・式は \label{名前} で名前を付け，\cref{名前} で参照する．
%% 「第2章」「3.1節」「図4.1」「式(3.1)」などの語と番号は自動で付く（手で書かない）．
本論文の構成は以下のとおりである．
\cref{chap:preliminaries}では，本研究で扱う問題の定義と，提案手法の基礎となる既存の解法について述べる．
\cref{chap:related-work}では，関連研究を整理し，本研究の位置づけを明らかにする．
\cref{chap:method}では，提案手法について述べる．
\cref{chap:experiments}では，数値実験の結果を示し，考察する．
\cref{chap:conclusion}では，本論文をまとめ，今後の課題を述べる．
